% Generated draft: requires mathematical and source review.
\documentclass[11pt]{article}
\usepackage[margin=25mm]{geometry}
\usepackage{fontspec}
\setmainfont{Latin Modern Roman}
\usepackage{amsmath,amssymb,mathtools,unicode-math}
\setmathfont{Latin Modern Math}
\usepackage{enumitem,needspace,xurl,hyperref}
\setlength{\parindent}{0pt}
\setlength{\parskip}{5pt}
\setlength{\emergencystretch}{3em}
\newcommand{\N}{\mathbb{N}}
\newcommand{\Z}{\mathbb{Z}}
\newcommand{\Q}{\mathbb{Q}}
\newcommand{\R}{\mathbb{R}}
\newcommand{\C}{\mathbb{C}}
\newcommand{\F}{\mathbb{F}}
\newcommand{\A}{\mathbb{A}}
\newcommand{\PP}{\mathbb{P}}
\newcommand{\E}{\mathbb{E}}
\newcommand{\eps}{\varepsilon}
\newcommand{\dd}{\,\mathrm{d}}
\begin{document}
\section*{UnsolvedMath ID: 800732}
\textbf{Source problem code:} catalogue entry
\subsection{Definitions and mathematical statement}
Let \(G\) be a group equipped with a topology. It is a \emph{semitopological group} if multiplication is separately continuous: for every \(a\in G\), both maps \(x\mapsto ax\) and \(x\mapsto xa\) are continuous. It is a \emph{quasitopological group} if, in addition, inversion \(x\mapsto x^{-1}\) is continuous. Joint continuity of multiplication is not required in either definition.

Use the usual two-sided interpretation of \emph{bounded}: for every neighborhood \(U\) of the identity there exist finite subsets \(F_L,F_R\subseteq G\) such that
\[
G=F_LU=UF_R,
\qquad F_LU=\{fu:f\in F_L,\ u\in U\},\quad UF_R=\{uf:u\in U,\ f\in F_R\}.
\]
The supplied excerpt does not explicitly define boundedness, so this interpretation requires confirmation against the source. The \emph{cellularity} of \(G\) is
\[
c(G)=\sup\bigl\{|\mathcal U|:\mathcal U\text{ is a pairwise disjoint family of nonempty open subsets of }G\bigr\}.
\]
Here \emph{countable cellularity} means \(c(G)\leq\aleph_0\), equivalently that every such family is at most countable.

The question attributed to Protasov asks whether every bounded semitopological group has countable cellularity:
\[
\forall G\quad
\bigl(G\text{ is a bounded semitopological group}\bigr)
\ \Longrightarrow\ c(G)\leq\aleph_0.
\]
The parenthetical wording also asks the same question for bounded quasitopological groups. These are recorded as two versions of the question, with the quasitopological version imposing continuous inversion. The supplied record does not specify separation axioms; the intended separation convention must therefore be checked before either version is regarded as a fully settled formulation.

For context, the supplied source reports that \(\sigma\)-compact topological groups have countable cellularity, citing [Tk], and that compact topological groups support a strictly positive Borel probability measure \(\mu\), meaning \(\mu(U)>0\) for every nonempty open subset \(U\). It further reports a result of T. Banakh and O. Ravsky, citing [BR] and also [BR2], that every bounded paratopological group has countable cellularity and, moreover, every cardinal of uncountable cofinality is a precaliber of that group. A paratopological group has jointly continuous multiplication, without requiring continuous inversion. A cardinal \(\kappa\) is a precaliber of a space if every indexed family of \(\kappa\) nonempty open sets has a subfamily indexed by a set of cardinality \(\kappa\) whose every finite subfamily has nonempty intersection. In contrast, the source reports, citing [Pr2], that for every infinite cardinal \(\tau\) there exists a left bounded left topological group with cellularity \(\tau\). These contextual assertions are supplied source reports, not independently verified results in this draft.
\subsection{Short English statement}
The question attributed to Protasov asks whether every bounded group with separately continuous multiplication has at most countably many pairwise disjoint nonempty open sets. It also asks the version in which inversion is continuous. The intended separation axioms and the source's exact convention for boundedness require confirmation. The registry labels the problem open; this draft does not independently verify that status.
\subsection{Sources}
\begin{itemize}
\item Local UnsolvedMath registry, supplied complete record: permanent ID \(800732\), source problem code \texttt{catalogue entry}; upstream revision \texttt{d371ee917ac9647fbb06cd6fc529deec6f3dda82}. The supplied status is \emph{open}.
\item Banakh and Choban, \emph{Some Open Problems in Topological Algebra} (2012), Problem 6, \url{https://arxiv.org/abs/1202.1619}. The registry lists Banakh and Choban as proposers and attributes the problem/context to Protasov. The supplied provenance records extraction by \texttt{arxiv\_latex} and access on \texttt{2026-10-16}; this is supplied dated provenance, not an independently verified visit or literature assessment.
\item The supplied excerpt cites [Tk] for the cellularity of \(\sigma\)-compact topological groups; [BR] and [BR2] for the results attributed to T. Banakh and O. Ravsky concerning bounded paratopological groups; and [Pr2] for left bounded left topological groups of prescribed infinite cellularity. Full bibliographic entries for these labels are not supplied. No separate citation is given for the assertion about strictly positive measures on compact topological groups.
\end{itemize}
\end{document}
